\section{Resultados e Discuss\~ao}
\label{sec:resultados}

\subsection{Precisão da evidência}
\label{sec:resultados-udv}

Os controles disponíveis indicam fragilidade no vínculo por similaridade semântica: nas 183 opiniões
de treino com citação confiável, a sentença de maior cosseno com a opinião completa é a da própria
citação em 112 (61,2\%). O intervalo de 95\% dessa proporção, por \textit{bootstrap} de audiências,
vai de 53,0\% a 68,8\%. Com os trechos entre aspas removidos da opinião, a sentença da citação é a
de maior cosseno em 12 de 142 consultas do treino e em 4 de 40 da validação. Esses números mostram
que a sentença escolhida por similaridade difere com frequência da passagem citada, mas não separam
o erro de outra sentença que também sustente a opinião. Essa separação cabe à validação humana, na
qual <resultado da validação humana: precisão estrita da citação literal e tolerante do nível
\texttt{semantic\_match\_high}, com intervalos de Wilson, critérios atingidos ou não, e concordância
do anotador consigo mesmo>.

\subsection{Efeito do perfil e dos trechos}
\label{sec:resultados-simulacao}

\input{tables/main-results}

Na validação, 81 das 102 opiniões ligadas geraram perguntas, de 51 atores em 21 audiências; as
demais não tinham três outros participantes na mesma audiência. Com $k = 3$, 5 e 10, o acerto da
condição 2 foi de 0,469, 0,481 e 0,481 no Gemma, com o empate resolvido em favor de $k = 5$ pela
probabilidade média da opção correta (0,467 contra 0,465), e de 0,494, 0,481 e 0,469 no Qwen, que
ficou com $k = 3$. No teste, 101 das 106 opiniões ligadas geraram perguntas, de 46 atores em 24
audiências (Tabela~\ref{tab:resultados}).

Nos dois modelos, a diferença de acerto entre as condições 1 e 0 foi positiva, de $+0{,}139$ no
Gemma e de $+0{,}119$ no Qwen, com intervalos que excluem o zero. O efeito dos trechos recuperados
não foi detectado, porque a diferença entre as condições 2 e 1, de $-0{,}020$ no Gemma e de
$+0{,}020$ no Qwen, tem intervalos que contêm o zero. Nas perguntas de validação, as condições 0, 1 e
2 tiveram acerto de 0,395, 0,543 e 0,481 no Gemma e de 0,358, 0,457 e 0,494 no Qwen, na mesma ordem
observada no teste em cada modelo, mas a condição 2 da validação é a do $k$ escolhido nessas mesmas
perguntas.

No teste, a soma das probabilidades das quatro letras ficou acima de 0,999 em todas as leituras das
condições 1 e 2 do Gemma e teve média de 0,929 e mínimo de 0,004 na condição 0. No Qwen, a média foi
de 0,999, 0,998 e 0,999 nas condições 0, 1 e 2, com mínimo de 0,969. Nas leituras da condição 0 do
Gemma com soma baixa, a resposta registrada resulta da renormalização de uma parcela pequena da
distribuição, e parte da diferença entre as condições 1 e 0 nesse modelo pode vir dessas leituras em
vez do conteúdo do perfil. No Qwen, cuja soma na condição 0 ficou acima de 0,99, a diferença teve
tamanho semelhante.

% Pendente: preencher a Figura fig:simulacao com uma pergunta de teste (itens do perfil, trechos recuperados, as quatro opções e as probabilidades nas condições 0 a 2), de um ator que não tenha entrado no teste preliminar; a regra de escolha da pergunta ainda precisa ser fixada antes da rodada final.

\subsection{Limitações}
\label{sec:limitacoes}

\textbf{Gabarito.} A opção correta é a opinião que a matéria atribui à pessoa, na redação da
matéria, e o acerto mede o reconhecimento dessa atribuição, que é uma seleção editorial do que foi
dito. Os distratores podem coincidir com posições que a pessoa também sustenta, sobretudo nas
condições que informam o lado dela, e essas perguntas admitem mais de uma resposta plausível. A
frequência dessas perguntas não foi medida.

\textbf{Informação comum às condições.} O nome e o cargo entram em todas as condições, e o cargo às
vezes registra a função da pessoa na própria audiência avaliada, como relator ou presidente da
comissão. O modelo pode associar ao nome posições conhecidas fora do conjunto. A condição 0 contém
esses efeitos e os do assunto e da redação das opções. As diferenças entre condições os cancelam
somente se eles se somarem ao efeito do material acrescentado, sem interagir com ele.

\textbf{Pré-treinamento.} As datas de corte do pré-treinamento do Gemma 4 31B, <data de corte>, e do
Qwen3.8 27B, <data de corte>, são <anteriores ou posteriores> ao período de teste. Se forem
posteriores, as matérias das audiências de teste podem ter sido vistas no pré-treinamento, com
efeito em todas as condições.

\textbf{Tamanho da amostra.} As perguntas de teste são 101, em 24 audiências, e um intervalo que
contém o zero indica que a diferença não foi detectada com essa amostra. Os dois modelos respondem
às mesmas perguntas, e por isso um resultado que se repete nos dois não constitui uma replicação em
amostra independente. O valor de $k$ é escolhido numa amostra de validação pequena e agrupada por audiência, em
que diferenças de poucas perguntas entre valores da grade podem refletir ruído.

\textbf{Desenho da avaliação.} Como as condições são encadeadas, a recuperação sem o perfil não é
avaliada, e o ganho do perfil sobre a simples recuperação das falas não é medido. A avaliação
também não separa assuntos já tratados nas falas de treino de assuntos novos.

\textbf{Perfis.} Cada modelo gerou os perfis uma única vez, com uma semente. A fidelidade dos perfis às falas não foi medida por anotação humana.

\textbf{Evidência da UDV.} A precisão da associação entre o nome citado na matéria e os turnos da
pessoa não foi medida. O corte de 0,45 entre \texttt{semantic\_match\_high} e
\texttt{semantic\_match\_weak} não foi validado como medida de confiança. Com negativos sorteados da
fala da própria pessoa, que é onde o codificador escolhe a evidência, a regra de pares daria 0,50. Nem
a construção nem o
verificador avaliam se a evidência sustenta a opinião, o que cabe à validação humana
(Seção~\ref{sec:setup-udv}).
